\subsection{A short injective resolution in nonnegative weights}

Write $\mathrm{gr}_{\geq0}S$ for the abelian category of
nonnegatively weight-graded $S=\Gamma_{\mathbb Z}(T)$-modules
and weight-preserving homomorphisms. Set
$S_p=S\otimes\mathbb Z_p$, $S_{\mathbb Q_p}=S\otimes\mathbb Q_p$,
and $E_p=S_{\mathbb Q_p}/S_p$, with quotient taken weight by
weight. The full integer-graded categories are used explicitly
where indicated below.

\paragraph{AT06 (Proved; completed coefficient injective dimension).}
The graded $S_p$-module $S_p$ has injective dimension at most
one in the full integer-graded module category. In
$\mathrm{gr}_{\geq0}S_p$, the displayed exact sequence
\[
 0\longrightarrow S_p\longrightarrow S_{\mathbb Q_p}
 \longrightarrow E_p\longrightarrow0
 \tag{AT.6}
\]
is an injective resolution.

\begin{proof}
Put $\overline S_p=S_p/pS_p=\Gamma_{\mathbb F_p}(T)$.
This is injective as a graded module over itself by
Nagpal--Snowden, Theorem~7.1(a), printed p.~39, whose
proof is on p.~41 of arXiv:1606.03431v2
\cite{IC-nagpal-snowden}. Its hypothesis holds because
the classical divided-power parameters at all positive
powers of $p$ vanish modulo $p$.

Let $M$ be any integer-graded $S_p$-module. The
two-term free resolution
$S_p\xrightarrow{p}S_p$ of $\overline S_p$ shows that
$M\otimes^{\mathbf L}_{S_p}\overline S_p$ has cohomology
only in degrees $-1,0$. Derived adjunction gives
\[
 R\operatorname{Hom}_{\mathrm{gr}S_p}(M,\overline S_p)
 =R\operatorname{Hom}_{\mathrm{gr}\overline S_p}
 (M\otimes^{\mathbf L}_{S_p}\overline S_p,\overline S_p).
\]
Injectivity of the target makes the right side have
cohomology only in degrees $0,1$. Thus
$\operatorname{Ext}^i(M,\overline S_p)=0$ for $i>1$.

Induction using
\[
 0\longrightarrow\overline S_p
 \xrightarrow{\ p^n\ }S_p/p^{n+1}S_p
 \longrightarrow S_p/p^nS_p\longrightarrow0
\]
gives $\operatorname{Ext}^i(M,S_p/p^nS_p)=0$ for $i>1$.
The same long exact sequence shows that the transition
maps
\[
 \operatorname{Ext}^1(M,S_p/p^{n+1}S_p)
 \longrightarrow\operatorname{Ext}^1(M,S_p/p^nS_p)
\]
are surjective.

Take a graded free resolution $P_\bullet\to M$. For
each cochain degree, the maps
\[
\operatorname{Hom}(P_i,S_p/p^{n+1}S_p)
\to\operatorname{Hom}(P_i,S_p/p^nS_p)
\]
are surjective:
each homogeneous free generator can be lifted separately.
Their inverse limit is
$\operatorname{Hom}(P_i,S_p)$, since the latter
coefficient module is complete in each weight.
For a countable inverse system of complexes with
surjective termwise transition maps, the exact sequence
\[
 0\to\varprojlim C_n\to\prod_n C_n
 \xrightarrow{\ 1-\mathrm{shift}\ }\prod_n C_n\to0
\]
and its long exact cohomology sequence give
\[
 0\to\varprojlim{}^1 H^{i-1}(C_n)
 \to H^i(\varprojlim C_n)
 \to\varprojlim H^i(C_n)\to0.
\]
Here surjectivity of the degree-one Ext transitions
makes the possible $\varprojlim^1$ term in degree two
zero; the other terms above degree one are already zero.
This proves $\operatorname{Ext}^i(M,S_p)=0$ for $i>1$.
It establishes the first assertion for arbitrary $M$,
without finite generation or flatness assumptions.

For the injective-resolution assertion, we spell out
the rational injectivity being used. If $V$ is a rational
vector space, the graded $S$-module $V[t]$, with
$T^{[a]}$ acting as $t^a/a!$, is injective in
$\mathrm{gr}_{\geq0}S$. Indeed for a nonnegative graded
$S$-module $L$, ordinary flat adjunction identifies
\[
 \operatorname{Hom}_{\mathrm{gr}S}(L,V[t])
 =\operatorname{Hom}_{\mathbb Q}
 \left(\varinjlim_{w\geq0}(L\otimes\mathbb Q)^w,V\right),
 \tag{AT.7}
\]
where transition maps are multiplication by $t=T$.
To check the formula, write a graded map in weight $w$
as $\ell\mapsto f_w(\ell)t^w$. Its $\mathbb Q[t]$-linearity
is exactly $f_{w+1}(t\ell)=f_w(\ell)$, the compatibility
for a map out of that colimit. Rationalization, taking
weight pieces, and this filtered colimit are exact;
$\operatorname{Hom}_{\mathbb Q}(-,V)$ is exact because
every short exact sequence of vector spaces splits.
This proves injectivity. The same argument over
$S_p$, after inverting $p$, uses $\mathbb Q_p$ in
place of $\mathbb Q$ and proves injectivity of
$S_{\mathbb Q_p}=\mathbb Q_p[t]$.

The nonnegative graded category computes the same Ext
as the full graded category for nonnegative source and
target, by the free resolutions specified before AT01.
Applying its long exact Ext sequence to (AT.6),
injectivity of $S_{\mathbb Q_p}$ and the first assertion
give $\operatorname{Ext}^1(M,E_p)=0$ for every
nonnegative $M$. Hence $E_p$ is injective in this
category too, proving (AT.6).
\end{proof}

\paragraph{AT07 (Proved; an explicit integral injective resolution).}
The object $S$ has injective dimension exactly one in
$\mathrm{gr}_{\geq0}S$. The following formula supplies
a length-one injective resolution.

Let $V=\prod_p\mathbb Q_p$, regarded as a rational
vector space, and identify $\mathbb A_f$ with its
subspace in $V$. Choose a $\mathbb Q$-linear retraction
$\rho:V\to\mathbb A_f$ of this inclusion. Put
$\lambda=\rho[t]:V[t]\to\mathbb A_f[t]$.
The graded embedding
$S_{\widehat{\mathbb Z}}\to V[t]$ sends
$bT^{[w]}$ to $(b/w!)t^w$.
Let
\[
 \overline{(-)}:V[t]=\prod_p S_{\mathbb Q_p}
       \longrightarrow\prod_p E_p
\]
be the coordinate quotient, where the products are graded
products. There is an exact sequence
\[
 0\longrightarrow S
 \xrightarrow{\ s\mapsto(s,s)\ }
 \mathbb Q[t]\oplus V[t]
 \xrightarrow{\ (a,b)\mapsto(\lambda(b)-a,\overline b)\ }
 \mathbb A_f[t]\oplus\prod_p E_p
 \longrightarrow0 .
 \tag{AT.8}
\]
Both displayed modules after $S$ are injective in
$\mathrm{gr}_{\geq0}S$.

\begin{proof}
A retraction $\rho$ exists by choosing a vector-space
complement of $\mathbb A_f$ in $V$. The formula makes
$\lambda$ an $S$-linear, weight-preserving map. On
$S_{\widehat{\mathbb Z}}$ it is the given inclusion into
$\mathbb A_f[t]$: each coefficient $b/w!$ belongs to
$\mathbb A_f$, and $\rho$ fixes that subspace.
All arrows in (AT.8) are therefore typed graded
$S$-module maps.

An element in the kernel has
$b\in S_{\widehat{\mathbb Z}}$ and $a=\lambda(b)=b$.
The intersection with $\mathbb Q[t]$ is precisely $S$,
weight by weight, by AT01. This proves exactness at
the middle term. To check surjectivity, let $(c,e)$
be homogeneous. Choose a homogeneous $b\in V[t]$
with $\overline b=e$, using the coordinate quotients.
By AT01 there are homogeneous
$a\in\mathbb Q[t]$ and $u\in S_{\widehat{\mathbb Z}}$
such that $\lambda(b)-c=a-u$. Then
$\overline{b+u}=e$ and
$\lambda(b+u)-a=c$. Thus $(a,b+u)$ is a preimage.
Finite sums of homogeneous preimages prove surjectivity
on the graded modules.

Formula (AT.7) proves injectivity of all the rational
coefficient modules occurring in (AT.8). Each $E_p$
is injective over $S_p$ by AT06 and remains injective
over $S$: restriction is right adjoint to the exact
flat scalar extension $-\otimes_S S_p$ in the
nonnegative graded categories. Products of injectives
are injective, since Hom into their product is the
product of exact Hom functors. Hence (AT.8) has the
claimed injective terms.

Finally $S/pS$ is a nonnegative graded module with
free resolution $0\to S\xrightarrow{p}S\to S/pS\to0$.
Its degree-preserving Ext into $S$ is
$\operatorname{Ext}^1(S/pS,S)=\mathbb Z/p$.
Thus $S$ is not injective, proving that its injective
dimension is exactly one.
\end{proof}

\paragraph{AT08 (Proved; failure of splitting for degreewise free abelian modules).}
The module $S$ is not injective in the exact category
of nonnegative graded $S$-modules whose individual
weight groups are free abelian, with exactness tested
weight by weight. More precisely, for
$1\leq j\leq N$ and $c=e_p\in\mathbb A_f\setminus\mathbb Q$,
shifting the explicit extension of AT03 gives a
nonsplit exact sequence
\[
 0\longrightarrow S\longrightarrow
 E_{c,N,j}\langle-j\rangle\longrightarrow
 I_N\langle-j\rangle\longrightarrow0
 \tag{AT.9}
\]
entirely within this exact category.

\begin{proof}
The source on the right has lowest weight $N-j\geq0$,
and every nonzero weight group is free of rank one.
The left term is likewise nonnegative and free of
rank one in every weight. In each weight the middle
term is an extension of these free abelian groups,
so that short exact sequence splits as abelian groups;
the middle term is therefore free abelian of rank at
most two in each weight. The weight-zero lower bound
also follows from the exact sequence. An $S$-linear
graded section of (AT.9), shifted back by $j$, would
be a section of AT03, where the explicit splitting
criterion is $c\in\mathbb Q$. Our choice of $c$ fails
that criterion. Thus (AT.9) is nonsplit despite all
the stated weightwise freeness properties.
\end{proof}

These are statements about explicitly specified graded
module morphisms. The resolution (AT.8) has not been
given a strict PD differential algebra structure or
used to construct an admissible common source. In
particular, its arbitrary vector-space retraction
$\rho$ is used only as a linear map.
